\section{Introduction}
\label{sec:intro}
End-to-end autonomous driving aspires to learn direct sensor-to-action policies~\cite{review1, review2, review3}, which hinges on effective visual representation learning~\cite{uniad, sparsedrive, zheng2024occworld, driveworld}. Propelled by recent breakthroughs in generative modeling, Driving World Models~(DWMs) are emerging as a promising paradigm for autonomous driving~\cite{review4, review5, review8}. By explicitly modeling the future scene evolution, DWMs provide a powerful foundation for forecasting the complex driving environments and downstream tasks.

\begin{figure}
    \centering
    \includegraphics[width=1.0\linewidth]{imgs/intro_v5.pdf}
    \caption{World models for end-to-end autonomous driving. (a)~Planning with future scenes generated by a driving world model. (b)~Planning with semantic representation extracted from a latent world model. (c)~WorldDrive bridges planning and driving world model via unifying vision and motion representation.}
    \label{fig:intro}
\end{figure}

Leveraging this predictive capability, the representation learned by DWMs holds immense potential for end-to-end planning. Current integration strategies generally follow two main paradigms. One approach leverages high-fidelity models to generate future scenes for downstream planners~\cite{drivewm, zeng2025fsdrive, zhao2025pwm, li2025imagidrive}. However, this process incurs prohibitive computational costs. To mitigate this overhead, a second paradigm operates entirely within a latent world model~\cite{law, world4drive, yang2025worldrft}. Although more efficient, this approach sacrifices interpretable scene simulation and limits explicit visual verification that is valuable for safety-critical decision-making. Crucially, we identify a systemic limitation pervading both strategies: representation misalignment and task disconnection. Scene simulators are typically optimized for perceptual reconstruction, emphasizing vision representation, whereas planners are trained in isolation for action regression to encode motion representation. This lack of a unified representation --- one that is shared and consistent across both scene generation and planning tasks --- prevents the planner from fully leveraging the generative driving world model’s learned dynamics and motion priors.

To bridge this gap, we introduce WorldDrive, a holistic framework designed to synergize end-to-end planning with scene generation via unifying vision and motion representation. The core philosophy of WorldDrive is representation unification: we posit that the latent features capable of generating the future~(scene generation) should be the same features used to decide the future~(planning). To this end, we first propose a Trajectory-aware Driving World Model~(TA-DWM). Unlike prior works that treat motion as a superficial condition, TA-DWM employs a multi-modal trajectory encoding scheme built on predefined trajectory anchors to construct a structured latent space where visual dynamics are intrinsically coupled with motion intentions. 

This design enables representation inheritance: the robust vision and motion encoders learned by the TA-DWM are directly transferred to initialize the downstream planner. This ensures that the planner operates in a mature and consistent latent space pre-aligned by the future scene generation task. Building upon this unified representation, we design a lightweight Multi-modal Planner. Leveraging these frozen encoders, the planner uses a query-centric cross-attention mechanism to efficiently fuse historical visual context with structured motion priors, generating diverse and high-quality trajectory candidates.

To harness the predictive foresight of the world model while avoiding the high latency of explicit video generation, we further introduce the Future-aware Rewarder~(FAR). Although TA-DWM is capable of synthesizing future videos, executing this generative process for every trajectory candidate is computationally prohibitive. Instead, the FAR employs a planning-oriented distillation mechanism to directly distill the future latents from the frozen world model. This approach enables WorldDrive to evaluate candidate trajectories based on their corresponding distilled future latent. This effectively aligns the planner’s selection with the world model’s learned dynamics while maintaining real-time inference speeds suitable for onboard deployment.

The main contributions of our work can be summarized as follows:


\begin{itemize}
    \item We introduce a novel trajectory-aware driving world model~(TA-DWM) that unifies the vision and motion representation. This design enables action-controllable future scene generation, producing plausible futures that are physically consistent with the conditioning trajectory.
    \item Leveraging the powerful representation learned by TA-DWM, we design a lightweight multi-modal planner that effectively fuses vision and motion cues. We also introduce a future-aware rewarder module that leverages the TA-DWM's foresight without the latency of explicit video generation, enabling real-time re-scoring of trajectory candidates at inference.
    \item We integrate these components into WorldDrive, a holistic framework that bridges the representational schism between visual simulation and trajectory planning. Its design enables two capabilities: high-fidelity, motion-consistent future scene generation, and multi-modal, real-time planning.
    \item We provide extensive evidence that WorldDrive achieves state-of-the-art WM-based end-to-end planning performance on NAVSIM, NAVSIM-v2, and nuScenes while concurrently achieving high-fidelity performance on conditional scene generation tasks.
\end{itemize}



